\section{Cash on hand as the state}\label{sec:cash}

\subsection{Saving is monotone, whatever the continuation}

\begin{theorem}[Monotone saving in cash on hand]\label{thm:topkis}
Let period utility be strictly concave on the positive half-line and let $W$ be \emph{any}
continuation value. If $m_1<m_2$ and $x_i$ is optimal at $m_i$ in \eqref{eq:step}, then
$x_1\le x_2$.
\formal{OLG.isOptimalSaving\_mono}
\end{theorem}

\begin{proof}[Proof sketch]
Suppose $x_2<x_1$. Both choices are feasible at both states, so adding the two optimality
inequalities,
\[
  u(m_1-x_1)+u(m_2-x_2)\ \ge\ u(m_1-x_2)+u(m_2-x_1),
\]
the continuation terms cancelling. Write $p=m_1-x_1$ and $q=m_2-x_2$. Since $x_2<x_1$ and
$m_1<m_2$ we have $p<m_1-x_2<q$, and $m_2-x_1=p+q-(m_1-x_2)$ also lies strictly between $p$ and
$q$. The two right-hand arguments are therefore strict convex combinations of $p$ and $q$ with
complementary weights, so strict concavity gives $u(m_1-x_2)+u(m_2-x_1)>u(p)+u(q)$, contradicting
the display.
\end{proof}

The proof uses nothing about $W$: it enters both optimality inequalities once and cancels. So a
value function made non-concave by a means test --- or by anything else --- cannot by itself make
saving fall with cash on hand. Note also what the hypotheses do \emph{not} include: no
differentiability, no interiority, no uniqueness of the optimum. The conclusion holds for arbitrary
optimal selections, which matters here, because at the switch point the optimum is genuinely
multi-valued.

\subsection{Consumption is not monotone}

What a non-concave continuation does produce is a jump in saving, and the jump is larger than the
increase in wealth that caused it.

\begin{theorem}[Consumption falls when wealth rises]\label{thm:cons}
There are a continuation value arising from a means-tested pension and two levels of cash on hand,
the second larger, at which optimal consumption is strictly smaller.
\formal{OLG.exists\_consumption\_decrease}
\end{theorem}

The witness is exact. Take logarithmic utility, a unit gross return, and a pension of one
withdrawn entirely above assets of one, so that next period's cash on hand is
\[
  N(x)=x+\mathbf{1}\{x\le1\},
\]
which jumps down by one as $x$ crosses the threshold, and let $W=\log N$. Then:

\begingroup\small
\begin{center}
\begin{tabular}{@{}lccc@{}}
\toprule
Cash on hand $m$ & optimal saving $x$ & consumption $m-x$ & next period's cash $N(x)$\\
\midrule
$6$ & $1$ & $5$ & $2$\\
$8$ & $4$ & $4$ & $4$\\
\bottomrule
\end{tabular}
\end{center}
\endgroup

At $m=6$ the household keeps the pension, saving the most it can without losing it; at $m=8$ the
pension is no longer worth protecting, so the household gives it up and saves the proceeds. Saving
rises from $1$ to $4$, as Theorem~\ref{thm:topkis} says it must, and since it rises by more than
wealth does, consumption falls from $5$ to $4$.

Both optima are verified rather than asserted. Because utility is logarithmic, comparing
$\log(m-x)+\log N(x)$ across $x$ is comparing the product $(m-x)N(x)$, which is a quadratic on each
side of the threshold. At $m=6$ the product is $-x^2+5x+6$ below the threshold, increasing there
and equal to $10$ at $x=1$, and $(6-x)x\le9$ above it. At $m=8$ it is $-x^2+7x+8\le14$ below and
$(8-x)x\le16$ above, attained at $x=4$. The comparisons are checked by the kernel
(\lean{isOptimalSaving\_six}, \lean{isOptimalSaving\_eight}).

\subsection{The two policies come apart}

The contrast is not an artefact of proof technique. Consumption is $m-x$ with both terms rising in
$m$, and when the continuation is non-concave the second can rise faster. Monotone consumption
therefore genuinely requires concavity of the continuation, while monotone saving does not; the
Lean development records the identity $c=m-x$ for exactly this purpose
(\lean{consumptionFnOf\_eq\_sub}) and derives consumption monotonicity only under concave slices.
So in the cash-on-hand state space the means test leaves the saving policy alone and breaks the
consumption policy, and the informal argument quoted in the introduction gets the object wrong as
well as the reason.
